\section{Kesimpulan \& Sintesis Akhir}

\begin{frame}{Ringkasan Perbaikan \& Validasi Teknis}
    \begin{itemize}
        \item \textbf{Integritas Matematis:} Reformulasi Hamiltonian dengan tanda negatif dan eliminasi \textit{double counting} telah menyelaraskan model dengan prinsip minimalisasi energi fisik.
        \item \textbf{Diagnostik Lanjutan:} Penggunaan \textit{Von Neumann Entropy} dan \textit{Gradient Variance} terbukti efektif dalam memvalidasi integritas state kuantum dan mendeteksi fenomena \textit{Barren Plateau}.
        \item \textbf{Akurasi Solusi:} Validasi melalui \textit{Brute Force} mengonfirmasi bahwa algoritma VQE mampu mencapai \textit{Ground State} portofolio dengan distribusi probabilitas yang akurat.
        \item \textbf{Optimalisasi GT:} Integrasi \textit{Game Theory Warm Start} secara signifikan meningkatkan kecepatan konvergensi dan stabilitas hasil investasi pada berbagai kondisi pasar.
    \end{itemize}
\end{frame}

\begin{frame}{Penutup}
    \centering
    \Huge{Sekian \& Terima Kasih}
    
    \vspace{1cm}
    \large{Diskusi \& Tanya Jawab}
\end{frame}
